\documentclass[a4paper,12pt]{article}

\usepackage[utf8]{inputenc}
\usepackage[T1]{fontenc}
\usepackage{amsmath}
\usepackage{booktabs}
\usepackage{graphicx}
\usepackage{geometry}
\geometry{margin=2.5cm}
\usepackage{xcolor}
\usepackage{listings}
\usepackage{hyperref}
\hypersetup{colorlinks=true, linkcolor=blue!50!black, urlcolor=blue!50!black}
\renewcommand{\tablename}{Tabela}
\renewcommand{\lstlistingname}{Código}

% ---------- Configuração do realce de sintaxe (C) ----------
\definecolor{codebg}{RGB}{247,247,247}
\definecolor{codekw}{RGB}{0,0,180}
\definecolor{codecomment}{RGB}{0,128,0}
\definecolor{codestring}{RGB}{163,21,21}
\definecolor{framecolor}{RGB}{200,200,200}

\lstdefinestyle{cstyle}{
    language=C,
    backgroundcolor=\color{codebg},
    basicstyle=\ttfamily\footnotesize,
    keywordstyle=\color{codekw}\bfseries,
    commentstyle=\color{codecomment}\itshape,
    stringstyle=\color{codestring},
    numberstyle=\tiny\color{gray},
    numbers=left,
    numbersep=6pt,
    frame=single,
    rulecolor=\color{framecolor},
    breaklines=true,
    breakatwhitespace=false,
    showstringspaces=false,
    tabsize=4,
    captionpos=b,
    columns=flexible,
    morekeywords={size_t},
    extendedchars=true,
    literate=
        {á}{{\'a}}1 {à}{{\`a}}1 {â}{{\^a}}1 {ã}{{\~a}}1
        {é}{{\'e}}1 {ê}{{\^e}}1 {í}{{\'i}}1
        {ó}{{\'o}}1 {ô}{{\^o}}1 {õ}{{\~o}}1
        {ú}{{\'u}}1 {ü}{{\"u}}1 {ç}{{\c c}}1
        {Á}{{\'A}}1 {À}{{\`A}}1 {Â}{{\^A}}1 {Ã}{{\~A}}1
        {É}{{\'E}}1 {Ê}{{\^E}}1 {Í}{{\'I}}1
        {Ó}{{\'O}}1 {Ô}{{\^O}}1 {Õ}{{\~O}}1
        {Ú}{{\'U}}1 {Ü}{{\"U}}1 {Ç}{{\c C}}1
}

\title{Tarefa 20: Equação do Calor 2D com Offload OpenMP para GPU}
\author{Israel Soares de Castro Filho \\ Matrícula: 20250070780}
\date{}

\begin{document}

\maketitle

\section{Introdução}

O objetivo da tarefa é executar na GPU o exemplo de transferência de calor
(\texttt{heat.c}) do tutorial de programação de GPUs com OpenMP, usando as
diretivas de movimentação de dados para que a cópia entre CPU e GPU não domine
a execução e a GPU não fique ociosa.

O programa resolve a equação do calor em uma grade $n \times n$ com um estêncil
de 5 pontos, avançando no tempo passo a passo. O resultado é validado contra a
solução exata conhecida (erro L2).

\section{Metodologia}

\textbf{Ambiente.} O código foi compilado com GCC, com OpenMP e offload para
NVIDIA (\texttt{nvptx}), e executado em uma GPU NVIDIA RTX 4070 Super local,
devido à fila de tarefas do NPAD. O perfilamento foi feito com o \texttt{nsys},
rastreando as chamadas CUDA.

\textbf{Estratégia de dados.} O ponto principal é manter os dois vetores
(\texttt{u} e \texttt{u\_tmp}) na GPU durante todo o laço de tempo:

\begin{lstlisting}[style=cstyle, caption={Diretivas de movimentação de dados}]
#pragma omp target enter data map(to: u[0:n*n], u_tmp[0:n*n])
for (int t = 0; t < nsteps; ++t) {
    solve(n, alpha, dx, dt, u, u_tmp);   // kernel na GPU
    tmp = u; u = u_tmp; u_tmp = tmp;     // troca de ponteiros
}
#pragma omp target exit data map(from: u[0:n*n])
\end{lstlisting}

Assim, os dados vão para a GPU uma vez, os 100 passos rodam sem nenhuma cópia
intermediária (a troca de ponteiros é feita no \textit{host}, sem copiar dados)
e só o resultado final volta para a CPU. O kernel usa
\texttt{\#pragma omp target} com \texttt{\#pragma omp loop collapse(2)}.

\textbf{Experimentos.} Foram testadas grades com $n \in \{500, 1000, 2000, 4000,
8000\}$, 100 passos de tempo e 3 repetições por tamanho. O programa mede o tempo
com o relógio do OpenMP e grava os resultados em um arquivo CSV. Cada perfil
do \texttt{nsys} cobre as 3 repetições (300 chamadas do kernel), e os tempos do CSV
foram medidos nessas mesmas execuções.

\section{Resultados}

A Tabela~\ref{tab:tempos} mostra as médias das 3 repetições. O tempo de
\textit{solve} é o laço de tempo na GPU (incluindo as cópias), e o tempo total
inclui também a inicialização e o cálculo do erro, feitos na CPU.

\begin{table}[h!]
\centering
\caption{Tempos médios (3 repetições, 100 passos).}
\label{tab:tempos}
\begin{tabular}{rrrrr}
\toprule
$n$ & Solve (s) & Total (s) & Fator vs. $n$ anterior & Erro L2 \\
\midrule
500  & 0,0445 & 0,0490 & --   & $8{,}10\times10^{-10}$ \\
1000 & 0,0906 & 0,1107 & 2,04 & $4{,}06\times10^{-10}$ \\
2000 & 0,2963 & 0,3764 & 3,27 & $2{,}02\times10^{-10}$ \\
4000 & 1,3071 & 1,6457 & 4,41 & $9{,}46\times10^{-11}$ \\
8000 & 5,3722 & 6,7985 & 4,11 & $5{,}14\times10^{-11}$ \\
\bottomrule
\end{tabular}
\end{table}

A Tabela~\ref{tab:perfil} resume o perfil do \texttt{nsys}. O kernel é o tempo
médio de uma chamada. Ida é o volume copiado da CPU para a GPU e volta é o da
GPU para a CPU, ambos por execução. A última coluna é o tempo das cópias (ida
mais volta) em relação ao tempo de \textit{solve}.

\begin{table}[h!]
\centering
\caption{Resumo do perfil (\texttt{nsys}).}
\label{tab:perfil}
\small
\begin{tabular}{rrrrrr}
\toprule
$n$ & Kernel (ms) & Kernel/Solve (\%) & Ida (MB) & Volta (MB) & Cópias/Solve (\%) \\
\midrule
500  & 0,398  & 89,5 & 4    & 2   & 0,8 \\
1000 & 0,784  & 86,5 & 16   & 8   & 1,7 \\
2000 & 2,785  & 94,0 & 64   & 32  & 2,0 \\
4000 & 12,75  & 97,6 & 256  & 128 & 1,7 \\
8000 & 52,70  & 98,1 & 1024 & 512 & 1,7 \\
\bottomrule
\end{tabular}
\end{table}

\section{Análise}

\textbf{Movimentação de dados.} Com \texttt{enter data} e \texttt{exit data}, os
vetores são copiados uma única vez por execução. As cópias ficam em torno de
2\% do tempo de \textit{solve} em todos os tamanhos, e o kernel ocupa de 87\% a
98\% dele. Ou seja, a GPU passa quase todo o tempo calculando, e não esperando
dados. Sem essas diretivas, cada passo copiaria os vetores nos dois sentidos, o
que daria, por estimativa (não medida), mais de 100 GB de tráfego para
$n = 8000$, contra cerca de 1,5 GB medidos aqui.

\textbf{Escalabilidade.} Para as grades menores, o tempo cresce menos que 4
vezes ao dobrar $n$ (fatores de 2,04 e 3,27), pois a GPU ainda não está
totalmente ocupada e os custos fixos (lançamento do kernel, alocação)
pesam. Para $n \geq 4000$ o fator se aproxima de 4, ou seja, o tempo passa a
crescer com o número de pontos ($n^2$).

\textbf{Tempo na CPU.} A diferença entre o tempo total e o de \textit{solve}
vem do trabalho feito na CPU (valor inicial e erro L2) e cresce com $n^2$: chega
a 1,43~s em $n = 8000$, cerca de 21\% do tempo total.

\textbf{Correção.} O erro L2 diminui à medida que a grade refina e é idêntico
entre as repetições, o que indica que o resultado na GPU está correto.

\section{Conclusão}

As diretivas \texttt{target enter data} e \texttt{target exit data} mantêm os
dados na GPU durante todo o laço de tempo. Com isso, as cópias ficaram em cerca
de 2\% do tempo de \textit{solve} e o kernel ocupa quase todo o restante, sem
GPU ociosa por falta de dados. Como a movimentação de dados deixou de ser o
problema, o próximo ganho de desempenho está na paralelização do próprio
kernel.

\newpage
\appendix

\section{Código-Fonte}

\lstinputlisting[style=cstyle, caption={tarefa20.c}]{../src/tarefa20.c}

\end{document}